\documentclass[11pt]{article}
\usepackage[T1]{fontenc}
\usepackage{amsmath,amssymb,amsthm}
\usepackage[margin=1in]{geometry}
\usepackage[colorlinks=true,urlcolor=blue]{hyperref}

\newtheorem*{theorem}{Theorem}

\title{An Explicit Counterexample to Erd\H{o}s Problem 280}
\author{}
\date{26 September 2026}

\begin{document}
\maketitle

\noindent\textbf{Provenance.}
This is an exposition and independent verification of a counterexample
already observed by Stijn Cambie and recorded by T.~F. Bloom in
\href{https://www.erdosproblems.com/280}{Erd\H{o}s Problem 280}.
It is not a claim of a new resolution of an open problem.

\medskip
\noindent\textbf{Question.}
Suppose $n_1<n_2<\cdots$ are positive integers and a residue class
$a_i\pmod{n_i}$ is chosen for every $i$. If for some $\varepsilon>0$
\begin{equation}\label{eq:growth}
 n_k>(1+\varepsilon)k\log k\qquad(k\geq 1),
\end{equation}
must the number of integers $0\leq m<n_k$ avoiding the first $k$
residue classes fail to be $o(k)$? The answer is no.

\begin{theorem}
Take
\[
 n_i=2^i,\qquad a_i=1+2^{i-1}\pmod{2^i}\qquad(i\geq1).
\]
Condition \eqref{eq:growth} holds with $\varepsilon=1$, but for every
$k\geq1$ one has
\[
 \#\bigl\{0\leq m<2^k:\;
       m\not\equiv a_i\pmod{2^i}\text{ for }1\leq i\leq k\bigr\}=1.
\]
Consequently this count is $o(k)$.
\end{theorem}

\begin{proof}
First, the moduli are strictly increasing. We check the numerical
growth condition without using an asymptotic estimate. For every
integer $k\geq2$, $\log k<k/2$: the function $x/2-\log x$ is positive
at $x=2$ and nondecreasing for $x\geq2$. Also $2^k\geq k^2$ for
$k\geq4$. Equality holds at $k=4$; the induction step follows from
$2k^2\geq(k+1)^2$ for $k\geq4$. Hence, for $k\geq4$,
\[
 2k\log k<k^2\leq2^k.
\]
The same strict inequality $2k\log k<2^k$ is immediate for
$k=1,2,3$. Thus \eqref{eq:growth} holds for $\varepsilon=1$.

Next, $m=1$ avoids every listed class, since
$a_i-1=2^{i-1}$ is not divisible by $2^i$.
If $0\leq m<2^k$ and $m\ne1$, let $t=v_2(m-1)$, the exponent of $2$
in the nonzero integer $m-1$ (this also makes sense when $m=0$).
Then $t\leq k-1$: this is clear for $m=0$, and otherwise
$1\leq m-1\leq2^k-2$. By the definition of $v_2$, we may write
$m-1=2^t u$ with $u$ odd, so
\[
 m\equiv1+2^t=a_{t+1}\pmod{2^{t+1}}.
\]
The index $t+1$ lies among $1,\ldots,k$. Therefore every $m$ in
$[0,2^k)$ except $1$ is covered by one of the first $k$ classes.
The uncovered count is exactly $1$, and $1/k\to0$.
\end{proof}

\paragraph{Why the construction works.}
The first $k$ congruence classes partition the residues modulo $2^k$
other than the residue $1$: the $i$th class consists exactly of those
integers for which $v_2(m-1)=i-1$. The one missing residue persists
while the moduli grow exponentially.

\paragraph{Source.}
P.~Erd\H{o}s and R.~L. Graham, \emph{Old and New Problems and Results
in Combinatorial Number Theory}, 1980, p.~29; problem statement and
Cambie's counterexample recorded by T.~F. Bloom,
\href{https://www.erdosproblems.com/280}{Erd\H{o}s Problem 280}.

\end{document}
